\documentclass[11pt]{article}

\usepackage[margin=1in]{geometry}
\usepackage[T1]{fontenc}
\usepackage[utf8]{inputenc}
\usepackage{amsmath,amssymb}
\usepackage{booktabs}
\usepackage{tabularx}
\usepackage{longtable}
\usepackage{enumitem}
\usepackage[hidelinks]{hyperref}

\newcommand{\code}[1]{\texttt{\detokenize{#1}}}
\newcommand{\rhoctl}{\rho}
\newcommand{\kappactl}{\kappa}
\newcommand{\Kbase}{K_{\mathrm{base}}}
\newcommand{\Keff}{K_{\mathrm{eff}}}
\newcommand{\Kref}{K_{\mathrm{ref}}}
\newcommand{\Kmax}{K_{\max}}
\newcommand{\reff}{r_{\mathrm{eff}}}
\newcommand{\ssafe}{s_{\mathrm{safe}}}

\setlength{\parskip}{0.35em}
\setlength{\parindent}{0pt}
\setlist[itemize]{leftmargin=1.4em,itemsep=2pt,topsep=2pt}
\setlist[enumerate]{leftmargin=1.6em,itemsep=2pt,topsep=2pt}
\sloppy

\title{\textbf{How Each Method Uses Public Ecology Parameters}\\
\large Algorithm-flow summary for the real-ecology benchmark}
\author{Distilled from server hidden-parameter audits}
\date{2026-07-14}

\begin{document}
\maketitle

\section{Purpose}

This note is a compact, algorithm-level explanation of where the current
real-ecology methods use quantities such as growth rate, carrying capacity, and
other table-derived ecological parameters. It deliberately avoids detailed code
walkthroughs. The goal is to remember the role of each algorithm and how
supposed-to-be-hidden quantities enter its learning, belief, planning, or solver
flow.

The server audit found that the current implementation is not mainly leaking a
private sidecar by accident. Instead, the current benchmark specification makes
several ecological quantities public. The methods then use those public
quantities in different ways.

\section{Current Known Setting vs Upcoming Hidden Setting}

The method-flow comparison in this note should be read against two different
experimental regimes.

\begin{center}
\begin{tabularx}{\linewidth}{p{0.24\linewidth}X}
\toprule
Regime & Meaning \\
\midrule
\textbf{Current run: public-demographics / parameter-known data-rich regime} &
This is the completed motivation-native run. It is a valid result, not a
discarded artifact. It shows that native mechanistic ecological solvers are very
strong when demographic/action tables are public. In this setting, the methods
can consume public controls such as \(\rhoctl\), \(\kappactl\), \(\Keff\),
\(\Kref\), action set-points, capacity effects, and K-derived safety quantities.
The result is therefore best interpreted as a data-rich or public-demographics
comparison, not as the intended hidden-parameter uncertainty test. Here,
``data-rich'' means the species' demographics/action effects are
well-characterized and public, not that the transition buffer is large. \\
\addlinespace
\textbf{Upcoming run: hidden-demographics / structure-unknown regime} &
This is the intended uncertainty test. ``Hidden demographics'' means the
methods may still know a population label if that is the chosen design, but they
must not use that label to look up \(r\), \(K\), caps, action effects,
\(\Kref\), or K-derived safety thresholds. The simulator/evaluator may still use
the real tables privately; the methods must infer the relevant growth/capacity
information from public trajectories. The structural family should also remain
private evaluator/simulator metadata, not method-facing metadata, so the run
tests hidden parameters plus model-form uncertainty rather than public table
lookup. \\
\bottomrule
\end{tabularx}
\end{center}

The structural axis should be kept separate from the parameter-visibility axis:

\begin{itemize}
  \item \textbf{Single-Ricker naive baseline:} \code{moor_native}. This is the
  headline naive ecological comparator because it remains locked to one Ricker
  form.
  \item \textbf{Closed-world form selection:} \code{plus_native}. This should be
  reported separately because it selects among the candidate families supplied
  to it; it is stronger than a naive single-form baseline.
  \item \textbf{Later open-world structural uncertainty:} a possible future
  \code{plus_native}-style experiment where candidate forms include distractors
  or the true family is sometimes absent. This should not be mixed into the first
  hidden-\(r,K\) repair, because it would confound the parameter-hiding fix with
  a new structural-model experiment.
\end{itemize}

\section{Current vs Hidden Fix by Method}

The table below is the most compact way to compare the completed known setting
with the upcoming hidden-\(r,K\) fix. It is written as a design target for the
server-side code change: the left column describes what the current run allowed,
and the right column describes what the corrected hidden setting should allow.
The schema fix is environment-level and applies to every method in the codebase;
the headline hidden-\(r,K\) run uses the five-method comparison
\code{refplan}, \code{bamcts}, \code{ogsrl}, \code{moor_native}, and
\code{plus_native}.

\begin{longtable}{p{0.17\linewidth}p{0.38\linewidth}p{0.38\linewidth}}
\toprule
Scope / method & Current run: public-demographics / parameter-known & Upcoming fix: hidden-demographics / structure-unknown \\
\midrule
\endfirsthead
\toprule
Scope / method & Current run: public-demographics / parameter-known & Upcoming fix: hidden-demographics / structure-unknown \\
\midrule
\endhead
General public schema &
Public dataset/belief/cache may contain \(\rhoctl\), \(\kappactl\), \(\Keff\),
\(\Kbase\), \(\Kref\), \(\Kmax\), action set-points, \(\Delta K\), \(\Delta N\),
K-derived \(\ssafe\), population table lookups, and family/config identity. &
Methods may consume only sanitized public data: population label as a categorical
token if chosen, \(o_t\), \(a_t\), action cost, reward as a label only,
\(o_{t+1}\), and done. All \(r/K\)-derived controls, table effects, K-derived
safety/normalization terms, and true family metadata stay private. \\

\code{refplan} &
Learns dynamics with \(\rhoctl\), \(\kappactl/\Kref\), and \(\Keff/\Kref\)
features; MPC planning carries public controls and uses \(\Kref\)/safety reward
terms. &
Fit dynamics and plan from sanitized belief/observation/action features only.
If action effects matter, infer them from data; do not pass exact \(\rhoctl\),
\(\kappactl\), \(\Keff\), \(\Kref\), or K-derived safety features. \\

\code{bamcts} &
Uses the same control-aware learned dynamics as \code{refplan}; tree keys and
rollouts condition on \(\rhoctl\) and \(\kappactl\). &
Tree state and rollouts must be indistinguishable for histories with the same
\((o,a)\) but different private \(r/K\). Learned dynamics cannot receive public
control aliases or exact table-derived action effects. \\

\code{ogsrl} &
Uses \(\rhoctl\), \(\kappactl/\Kref\), \(\Keff/\Kref\), \(\ssafe\), and unsafe
features in dynamics, actor features, guardian features, rollouts, and safety
checks. &
Remove \(r/K\)-derived features from dynamics, actor, guardian, rollout, and
safety inputs. Safety-aware behavior must learn risk from public trajectories or
use a public threshold that is not derived from hidden \(\Kbase\). \\

\code{mopo} &
Uses the same learned-dynamics public-control channel as the other general
planners; planning reward uses \(\Kref\) and safety configuration. &
Use sanitized features only. Model uncertainty should reflect uncertainty about
hidden action effects and capacity rather than conditioning on exact public
controls. \\

\code{delphic} &
Does not build an ecological transition model, but receives shared belief
features containing \(\rhoctl\), \(\kappactl/\Kref\), \(\Keff/\Kref\), and
safety terms. &
Compatible-world and Q features must be built from sanitized public beliefs
only. Remove hidden-parameter aliases from the shared feature vector so
\code{delphic} cannot inherit the leak indirectly. \\

adapted \code{moor} &
Uses exact action set-points, \(\Delta K\), \(\Delta N\), and species capacity
structure in a mechanistic Ricker proposal; mainly fits a K-like scale in
set-point mode. &
Mechanistic proposal cannot call table-derived exact action effects or species
capacity. It must fit required Ricker parameters/effects from public data, or be
reported only in the public-demographics comparison arm. \\

adapted \code{plus} &
Uses exact table-derived action effects and species capacity ranges inside a
candidate mechanistic bank; candidate axis is mostly \(K\) because growth
set-points are public. &
Candidate parameters and priors must be derived from sanitized public data, not
from species/action tables. If retained, report separately from the naive
headline if it performs closed-world model selection. \\

\code{moor_native} &
Fits \(\hat K\), but reads exact action growth set-points, \(\Delta K\),
\(\Delta N\), costs, \(\Kref\), and \(\ssafe\) to build its native Ricker solver. &
Headline naive baseline. It should fit/estimate Ricker parameters and action
effects from sanitized public data, then solve a single-Ricker model. It must not
read species/action tables or K-derived safety/normalization values. \\

\code{plus_native} &
Builds one native solver per candidate family using exact table-derived
\(r/K\), action effects, reward, and safety structure; updates mainly the family
posterior. &
Separate closed-world model-selection baseline. It may keep a family posterior,
but all parameter candidates/priors must be fit from sanitized public data. Do
not redesign open-world PLUS in the first hidden-\(r,K\) repair. \\
\bottomrule
\end{longtable}

\section{Shared Benchmark Information Flow}

The simulator privately owns the real ecology tables. These tables contain, for
each population and action, quantities such as:
\[
  \lambda(a,p),\quad r(a,p),\quad \Delta K(a,p),\quad \Delta N(a,p),
  \quad \Kbase(p),\quad \Kmax(p).
\]

In the current public-information design, the methods do not just see
observations and actions. They also receive or derive public controls such as:
\[
  \rhoctl,\quad \kappactl,\quad \Keff,\quad \Kref,\quad \ssafe,
\]
where \(\rhoctl\) is the action growth set-point, \(\kappactl\) is cumulative
capacity investment, \(\Keff\) is effective carrying capacity, \(\Kref\) is the
reward normalizer, and \(\ssafe\) is a safety threshold.

\begin{center}
\begin{tabularx}{\linewidth}{p{0.23\linewidth}X}
\toprule
Quantity & Current meaning in the benchmark \\
\midrule
\(\rhoctl\), \(\reff\) & Growth-rate control. In the real set-point setting, the action table effectively tells the method the growth effect of the chosen action. \\
\(\kappactl\), \(\Keff\) & Capacity-control state. These tell the method how capacity interventions accumulate and what effective capacity is being used. \\
\(\Kbase\), \(\Kref\), \(\Kmax\) & Species-scale capacity quantities. These expose or normalize by the population's capacity scale. \\
\(\Delta K\), \(\Delta N\) & Exact action effects for restoration and translocation. \\
\(\ssafe\) & Safety threshold, currently derived from capacity, so it can indirectly reveal \(\Kbase\). \\
Population identity & Public label. If it can be used to look up species tables, it becomes a demographic side channel. \\
Family identity / \code{cfg.kind} & The true dynamics family or configuration identity. It should be treated carefully because it can leak structural information. \\
\bottomrule
\end{tabularx}
\end{center}

\section{The Big Distinction}

The important distinction is not simply whether a method sees \(r\) or \(K\).
The important distinction is how the method uses that information.

\begin{center}
\begin{tabularx}{\linewidth}{p{0.28\linewidth}X}
\toprule
Method type & How public ecology parameters are used \\
\midrule
General learned methods & Use \(r/K\)-derived quantities as features or rollout context for learned dynamics and policies. They still learn an approximate model from data. \\
Mechanistic adapted methods & Use exact table-derived action/species structure inside a mechanistic proposal, then plan with particle MPC. \\
Native ecological solvers & Build discrete ecological transition/reward solvers from exact public ecological structure. This is the strongest use of the tables. \\
\bottomrule
\end{tabularx}
\end{center}

So the current asymmetry is:
\[
\text{general methods use public controls as features}
\quad\text{whereas}\quad
\text{ecological solvers use them as model ingredients.}
\]

\section{Method-by-Method Summary}

\begin{longtable}{p{0.17\linewidth}p{0.29\linewidth}p{0.30\linewidth}p{0.16\linewidth}}
\toprule
Method & Algorithm context & How \(r/K\)-like quantities enter & Main consequence \\
\midrule
\endfirsthead
\toprule
Method & Algorithm context & How \(r/K\)-like quantities enter & Main consequence \\
\midrule
\endhead
\code{refplan} &
General offline model-based planner. It learns an ensemble of dynamics models and plans under a posterior over those models. &
The learned dynamics features include \(\rhoctl\), \(\kappactl/\Kref\), and \(\Keff/\Kref\). Planning rolls particles forward while carrying public controls. &
Uses public \(r/K\) features, but still learns dynamics from data. \\

\code{bamcts} &
Bayes-adaptive tree-search style planner. It learns a dynamics ensemble, then searches over actions in a belief/model tree. &
The dynamics model uses the same public control features. The search tree also buckets or simulates with \(\rhoctl\) and \(\kappactl\). Reward and safety use \(\Kref\) and \(\ssafe\). &
Tree search conditions directly on public control state. \\

\code{ogsrl} &
Offline policy-learning method with support and safety guards. It has a learned dynamics model, a KNN out-of-distribution guardian, a policy head, and safety checks. &
Uses \(\rhoctl\), \(\kappactl/\Kref\), and \(\Keff/\Kref\) in dynamics fitting, guardian features, policy features, rollout simulation, and unsafe/safety features. &
Broadest \(r/K\)-feature exposure among the general methods. \\

\code{mopo} &
General model-based offline RL baseline. It learns a bootstrap dynamics ensemble and penalizes uncertain model rollouts. &
Uses public controls in the learned dynamics model and particle-MPC rollouts, similar to \code{refplan}. &
Same public-control feature channel as the other learned planners. \\

\code{delphic} &
Compatible-world / conservative Q-learning style method. It does not build an explicit ecological transition model. &
Receives shared belief features. Those belief features already include \(\rhoctl\), \(\kappactl/\Kref\), \(\Keff/\Kref\), and safety-related terms. &
Indirect \(r/K\) exposure through the shared feature vector. \\

adapted \code{moor} &
Mechanistic Ricker baseline adapted to the shared particle-MPC infrastructure. It fits a simple Ricker-style model and plans with MPC. &
In the real set-point setting, action growth is already known through public \(\rhoctl\). The method mainly fits a \(K\)-like scale, while using exact action set-points, capacity increments, and reward/safety structure. &
Single-Ricker and misspecified off-Ricker, but not table-blind. \\

adapted \code{plus} &
Mechanistic PLUS-style particle-MPC adaptation. It keeps a bank of candidate mechanistic models and updates candidate weights. &
Candidate models use exact action set-points and action effects. In set-point mode, the candidate axis is mostly \(K\), because action growth is already public. &
Mechanistic planner advantage, though still using particle MPC rather than native discrete solving. \\

\code{moor_native} &
Native discrete ecological solver. It is the cleaner naive ecological baseline because it commits to one Ricker form. &
It fits a scalar \(\hat K\), but still reads exact action growth set-points, \(\Delta K\), \(\Delta N\), action costs, \(\Kref\), and safety threshold. It does not learn action growth effects from \((o,a,o')\). &
Naive in functional form, but still uses public demographic/action tables. \\

\code{plus_native} &
Native discrete ecological model-selection solver. It builds one solver for each candidate family and updates a posterior over the family. &
It builds candidate solvers from exact table-derived \(r\), \(K\), \(\Delta K\), \(\Delta N\), reward, and safety structure. It does not estimate \(r/K\) from public transition data; it mainly updates the family posterior. &
Closed-world model-identifying solver, not a naive baseline. \\
\bottomrule
\end{longtable}

\section{General Learned Methods}

\subsection{\code{refplan}}

\textbf{Algorithm context.} \code{refplan} is a general offline model-based RL
method. In this benchmark it learns an ensemble of dynamics models from the
offline dataset. At decision time, it treats the ensemble as a posterior over
possible dynamics models and plans by particle MPC.

\textbf{Flow.}
\[
\text{public dataset}
\rightarrow
\text{belief cache}
\rightarrow
\text{learn ensemble dynamics}
\rightarrow
\text{posterior over ensemble members}
\rightarrow
\text{MPC planning}.
\]

\textbf{Where \(r/K\)-like information enters.} The dynamics model is fitted
using features that include \(\rhoctl\), \(\kappactl/\Kref\), and
\(\Keff/\Kref\). During planning, the MPC rollout carries \(\rhoctl\) and
\(\kappactl\) forward and scores reward using \(\Kref\) and safety settings.

\textbf{Interpretation.} \code{refplan} does use public ecology controls, but it
does not get a ready-made ecological equation in the same way the native solvers
do. It still has to learn a dynamics model.

\subsection{\code{bamcts}}

\textbf{Algorithm context.} \code{bamcts} is a Bayes-adaptive planning method.
In this benchmark it uses a learned ensemble dynamics model and performs a
tree-search style action selection over belief/model uncertainty.

\textbf{Flow.}
\[
\text{public dataset}
\rightarrow
\text{learn ensemble dynamics}
\rightarrow
\text{tree search over actions and model samples}
\rightarrow
\text{execute best action}.
\]

\textbf{Where \(r/K\)-like information enters.} The learned dynamics uses the
same public-control features as \code{refplan}. The search state also includes
or buckets \(\rhoctl\) and \(\kappactl\), and simulated rollouts advance these
controls.

\textbf{Interpretation.} \code{bamcts} is explicitly conditioned on public
growth/capacity controls during search. It is not blind to \(r/K\), but the
transition model is still learned rather than supplied as a correct ecological
model.

\subsection{\code{ogsrl}}

\textbf{Algorithm context.} \code{ogsrl} is an offline policy method with support
and safety guards. It learns a policy but also uses a dynamics model, a KNN
out-of-distribution guardian, safety critics/features, and a feasibility mask at
deployment.

\textbf{Flow.}
\[
\text{public belief features}
\rightarrow
\text{learn dynamics and support guardian}
\rightarrow
\text{train guarded policy}
\rightarrow
\text{roll out / check action support and safety}
\rightarrow
\text{act}.
\]

\textbf{Where \(r/K\)-like information enters.} \code{ogsrl} uses
\(\rhoctl\), \(\kappactl/\Kref\), and \(\Keff/\Kref\) in several places:
dynamics fitting, policy features, KNN guardian features, rollout simulation,
and safety/unsafe features. It also uses \(\ssafe\), which can leak capacity if
derived from \(\Kbase\).

\textbf{Interpretation.} Among the general methods, \code{ogsrl} has the widest
public \(r/K\)-derived feature exposure. But its use is still feature-based and
learned; it is not simply solving with the exact ecological tables.

\subsection{\code{mopo}}

\textbf{Algorithm context.} \code{mopo} is a model-based offline RL method. It
learns a dynamics ensemble and penalizes planning in regions where the learned
model is uncertain.

\textbf{Flow.}
\[
\text{offline data}
\rightarrow
\text{bootstrap dynamics ensemble}
\rightarrow
\text{uncertainty-penalized MPC planning}.
\]

\textbf{Where \(r/K\)-like information enters.} It uses the same learned-dynamics
feature channel as \code{refplan}: \(\rhoctl\), \(\kappactl/\Kref\),
\(\Keff/\Kref\), plus reward/safety quantities for planning.

\textbf{Interpretation.} \code{mopo} receives public ecology controls as model
features, but the dynamics remain learned and approximate.

\subsection{\code{delphic}}

\textbf{Algorithm context.} \code{delphic} is a compatible-world / conservative
Q-learning style method. It does not primarily learn an ecological transition
equation. Instead, it creates multiple compatible latent explanations and
penalizes actions whose values vary across those explanations.

\textbf{Flow.}
\[
\text{belief features}
\rightarrow
\text{compatible worlds}
\rightarrow
\text{world-specific value estimates}
\rightarrow
\text{conservative Q policy}.
\]

\textbf{Where \(r/K\)-like information enters.} It receives the shared belief
features. Since those features include \(\rhoctl\), \(\kappactl/\Kref\), and
\(\Keff/\Kref\), \code{delphic} indirectly receives public \(r/K\)-derived
information.

\textbf{Interpretation.} Its exposure is indirect. It does not use tables to
construct a mechanistic ecological planner, but it still sees public control
features.

\section{Adapted Ecological Methods}

\subsection{Adapted \code{moor}}

\textbf{Algorithm context.} Adapted \code{moor} is a mechanistic ecological
baseline fitted into the shared particle-MPC framework. It uses a single Ricker
form and plans with particle MPC.

\textbf{Flow.}
\[
\text{public belief means and actions}
\rightarrow
\text{fit Ricker-style K}
\rightarrow
\text{mechanistic proposal}
\rightarrow
\text{particle-MPC planning}.
\]

\textbf{Where \(r/K\)-like information enters.} In the real set-point setting,
the action growth effect is already public through \(\rhoctl\). Therefore the
method mainly fits a \(K\)-like scalar while relying on exact public action
growth set-points, capacity increments, translocation magnitude, and reward
structure.

\textbf{Interpretation.} It is naive in the sense of using a single Ricker form,
but it is not blind to the action/species table.

\subsection{Adapted \code{plus}}

\textbf{Algorithm context.} Adapted \code{plus} is a PLUS-style mechanistic
method adapted to the shared particle-MPC infrastructure. It keeps a bank of
candidate mechanistic models and updates weights/posteriors over those
candidates.

\textbf{Flow.}
\[
\text{candidate mechanistic models}
\rightarrow
\text{candidate belief banks}
\rightarrow
\text{posterior update from observations}
\rightarrow
\text{posterior-weighted MPC planning}.
\]

\textbf{Where \(r/K\)-like information enters.} Candidate proposals use exact
public action set-points, capacity increments, translocation magnitudes, and
species-scale capacity ranges. Because action growth is already public, the
candidate axis becomes mostly \(K\)-like rather than action-growth-like.

\textbf{Interpretation.} This is a structured mechanistic method with public
table access. It is not the same as a method that must infer all action effects
from \((o,a,o')\).

\section{Native Ecological Solvers}

\subsection{\code{moor_native}}

\textbf{Algorithm context.} \code{moor_native} is the native discrete ecological
baseline intended to be the clean naive comparator. It uses a native discrete
belief/filter and tabular value-solving machinery, while committing to a single
Ricker form.

\textbf{Flow.}
\[
\text{public observations/actions}
\rightarrow
\text{fit } \hat K
\rightarrow
\text{build native Ricker transition/reward table}
\rightarrow
\text{discrete belief update}
\rightarrow
\text{tabular planning}.
\]

\textbf{Where \(r/K\)-like information enters.} It fits \(\hat K\), but it still
uses exact action growth set-points, exact \(\Delta K\), exact \(\Delta N\),
action costs, \(\Kref\), and \(\ssafe\). It does not estimate the action growth
effects from public transition data.

\textbf{Interpretation.} This is the best ``naive'' ecological baseline only if
naive means single-form Ricker. It is not naive in the sense of not knowing the
public action/species mechanisms.

\subsection{\code{plus_native}}

\textbf{Algorithm context.} \code{plus_native} is a native PLUS-style solver. It
builds separate native solvers for candidate ecological families, then updates a
posterior over which family best explains the observations.

\textbf{Flow.}
\[
\text{exact public ecology tables}
\rightarrow
\text{one native solver per candidate family}
\rightarrow
\text{posterior over family}
\rightarrow
\text{posterior-weighted native action values}.
\]

\textbf{Where \(r/K\)-like information enters.} It builds candidate solvers from
exact table-derived action growth, capacity, translocation, reward, safety, and
species-scale information. It does not estimate \(r\) or \(K\) from public
transition data. It mainly estimates or updates the family posterior.

\textbf{Interpretation.} \code{plus_native} should not be described as the naive
baseline. It is better described as a closed-world model-selection ecological
solver with exact public demographic/action structure.

\section{Practical Takeaways}

\begin{enumerate}
  \item The problem is not just that \(r/K\)-like values are visible. They are
  actively used in method flows.
  \item General methods use them mostly as learned-model features and rollout
  context.
  \item Mechanistic and native ecological methods use them as stronger
  transition-model ingredients.
  \item \code{moor_native} partly estimates \(K\), but still reads exact action
  effects and safety/reward structure.
  \item \code{plus_native} does not estimate \(r/K\); it updates model-family
  beliefs while using exact public table structure.
  \item If the intended experiment is hidden-demographics or parameter
  uncertainty, the public/private boundary must hide action growth, capacity,
  translocation magnitude, \(\Kref\), \(\ssafe\), and table-derived aliases from
  all method-facing paths.
\end{enumerate}

\section{One-Sentence Summary}

In the current benchmark, general methods use public \(r/K\)-derived controls as
features for learned dynamics, while mechanistic and native ecological methods
use the same public information as ingredients of much stronger ecological
planning models; therefore the completed run is best described as a valid
public-demographics / parameter-known data-rich regime, while the upcoming
hidden-demographics run is the intended hidden-\(r,K\) and structure-unknown
uncertainty test.

\end{document}
